\documentclass[a4paper,12pt]{article}
\usepackage[utf8]{inputenc}
\usepackage[english, russian]{babel}
\usepackage{amsmath, amssymb, amsthm}
\usepackage{enumitem}%оформление списков
\usepackage{varwidth}
\usepackage{graphicx}
\usepackage{float}
\usepackage{indentfirst}
\usepackage{url}
\usepackage{seqsplit}
\usepackage{pdflscape}
\usepackage{tabularx}
\usepackage{xcolor}
\usepackage{listings}
\usepackage{mathtools}
\usepackage{hyperref}
\usepackage[backend=biber, style=numeric, sorting=none]{biblatex}
\usepackage{titlesec}
\usepackage{appendix}
\usepackage{alltt}

\addbibresource{my_libs.bib}


\usepackage[a4paper, top=20mm, left=30mm, right=20mm, bottom=20mm]{geometry}
\newcommand{\dim}{\mathrm{dim}}
\newcommand{\lin}{\mathrm{lin}}
\newcommand{\diag}{\mathrm{diag}}
\newcommand{\R}{\mathbb{R}}
\newcommand{\Rn}{\mathbb{R^n}}

\lstset{
    language=C++,
    basicstyle=\ttfamily\small,
    numbers=left,
    stepnumber=1,
    breaklines=true,
    breakatwhitespace=true,
    keepspaces=true,
    tabsize=4,
    showspaces=false,
    showstringspaces=false,
    extendedchars=\true,
}

\begin{document}
\appendix
\section{Исходный код}
\begin{lstlisting}
// Файл: ./src/main.cpp

#include "cli/CommandLineParser.h"
#include "../include/processing/Processing.h"
#include <iostream>

int main(int argc, char* argv[]) {
//  РАЗБОР АРГУМЕНТОВ КОМАНДНОЙ СТРОКИ
    ProgramOptions opts = CommandLineParser::parse(argc, argv);

    if (opts.showHelp) {
        CommandLineParser::printHelp(argv[0]);
        return 0;
    }

    if (!opts.partial) {
        std::cout << "Режим: полный перебор всех вершин" << std::endl;
    }

    if (opts.batchMode && !opts.inputDir.empty()) {
//  ПАКЕТНАЯ ОБРАБОТКА
        return Processing::processBatch(opts);
    }

    if (opts.loadFromFile && !opts.inputFile.empty()) {
//  ОБРАБОТКА ОДИНОЧНОГО ФАЙЛА
        return Processing::processSingleFile(opts.inputFile, opts);
    }

// СТАНДАРТНОЕ ПОСТРОЕНИЕ (СИЛЬВЕСТР)
    return Processing::processSilvester(opts);
}

=========================================================
// Файл: ./include/basis/SimplexBasis.h

#pragma once
#include <Eigen/Dense>
#include <vector>

class SimplexBasis {
   public:
    explicit SimplexBasis(const std::vector<Eigen::RowVectorXd>& vertices);
    [[nodiscard]] int    size() const;
    [[nodiscard]] double value(int i, const Eigen::RowVectorXd& x) const;
    static bool          verifyInterpolation(const SimplexBasis& basis, double tolerance);
    [[nodiscard]] const std::vector<Eigen::RowVectorXd>& getVertices() const {
        return vertices_;
    }

   private:
    std::vector<Eigen::RowVectorXd> vertices_;
    int                             dim_;
    Eigen::MatrixXd                 Ainv_;

   public:
    [[nodiscard]] Eigen::MatrixXd getAinv() const {
        return Ainv_;
    }
};

=========================================================
// Файл: ./include/cli/CommandLineParser.h

#pragma once
#include <string>
#include <cstdint>

struct ProgramOptions {
// Режим перебора
    bool    partial     = false;
    int64_t startVertex = 0;
    int64_t endVertex   = 0;
    int     dimension   = 7;

// Параметры выполнения
    int numThreads = 0;  // 0 = автоматически (hardware_concurrency)
    int64_t logInterval = 10'000'000;  // интервал логирования (количество вершин)

// Режим загрузки данных
    bool        loadFromFile = false;
    bool        batchMode    = false;
    std::string inputFile;
    std::string inputDir;
    enum class FileType { HADAMARD } fileType = FileType::HADAMARD;

// Помощь
    bool showHelp = false;
};

class CommandLineParser {
   public:
    static ProgramOptions parse(int argc, char* argv[]);
    static void           printHelp(const char* programName);

   private:
    static void parseRange(ProgramOptions& opts, int& i, char* argv[]);
    static void parseDimension(ProgramOptions& opts, int& i, char* argv[]);
    static void parseThreads(ProgramOptions& opts, int& i, char* argv[]);
    static void parseLogInterval(ProgramOptions& opts, int& i, char* argv[]);
};

=========================================================
// Файл: ./include/hadamard_utils/HadamardMatrixIterator.h

#pragma once
#include <string>
#include <vector>
#include <fstream>
#include <Eigen/Dense>

class HadamardMatrixIterator {
   public:
    explicit HadamardMatrixIterator(const std::string& filename);

    void            reset();
    bool            hasNext() const;
    Eigen::MatrixXd next();

   private:
    std::string   filename_;
    std::ifstream file_;
    int           currentOrder_;

    static bool isCommentLine(const std::string& line);
    static bool isSymbolicLine(const std::string& line);
    static bool isNumericLine(const std::string& line);

    static std::vector<double> parseLine(const std::string& line);
    Eigen::MatrixXd            readNextMatrix();
};

=========================================================
// Файл: ./include/hadamard_utils/HadamardUtils.h

#pragma once
#include <Eigen/Dense>
#include <vector>

class HadamardUtils {
   public:
    static std::vector<Eigen::RowVectorXd> getVerticesWithLogging(const Eigen::MatrixXd& H,
                                                                  int                    dim);
    static bool isHadamardMatrix(const Eigen::MatrixXd& H, double tolerance = 1e-10);

   private:
// Преобразование матрицы Адамара в (0/1)-матрицу D
    static Eigen::MatrixXd hadamardTo01Matrix(const Eigen::MatrixXd& H);

// Построение вершин симплекса из (0/1)-матрицы D
    static std::vector<Eigen::RowVectorXd> buildVerticesFrom01Matrix(const Eigen::MatrixXd& D);
};

=========================================================
// Файл: ./include/processing/Processing.h

#pragma once
#include <string>
#include <Eigen/Dense>

struct ProgramOptions;

class Processing {
    static int    processHadamardMatrixFromFile(const std::string&    filename,
                                                const ProgramOptions& opts, double& globalMin,
                                                int&                globalMinMatrixIdx,
                                                Eigen::RowVectorXd& globalMinPoint);
    static double processHadamardMatrix(const Eigen::MatrixXd& H, const ProgramOptions& opts,
                                        int matrixCounter);

   public:
    static int processBatch(const ProgramOptions& opts);
    static int processSingleFile(const std::string& filename, const ProgramOptions& opts);
    static int processSilvester(const ProgramOptions& opts);
};

=========================================================
// Файл: ./include/upper_bound_calculator/FullEnumUpperBoundCalculator.h

#pragma once
#include "upper_bound_calculator/UpperBoundCalculatorInterface.h"
#include "upper_bound_calculator/LebesgueFunction.h"
#include <atomic>
#include <Eigen/Dense>

class FullEnumUpperBoundCalculator : public UpperBoundCalculatorInterface {
   public:
    FullEnumUpperBoundCalculator(const LebesgueFunction& func, int dim);

    FullEnumUpperBoundCalculator(const LebesgueFunction& func, int dim, int64_t startVertex,
                                 int64_t endVertex);

// Запуск перебора. Возвращает максимальное значение константы Лебега
    double calculate() override;

    [[nodiscard]] Eigen::RowVectorXd getMaxPoint() const override {
        return maxPoint_;
    }

   private:
    const LebesgueFunction& func_;
    int                     dim_;
    Eigen::RowVectorXd      maxPoint_;
    std::atomic<double>     globalMax_;
    std::atomic<int64_t>    processedVertices_;
    std::atomic<int64_t>    lastLogged_;
    std::atomic<int64_t>    currentBestVertex_;

// Параметры частичного перебора
    int64_t startVertex_;
    int64_t endVertex_;
    bool    partial_;
};

=========================================================
// Файл: ./include/upper_bound_calculator/LebesgueFunction.h

#pragma once
#include "basis/SimplexBasis.h"

class LebesgueFunction {
   public:
    explicit LebesgueFunction(const SimplexBasis& basis);
    double operator()(const Eigen::RowVectorXd& x) const;

   private:
    const SimplexBasis& basis_;
};

=========================================================
// Файл: ./include/upper_bound_calculator/UpperBoundCalculatorInterface.h

#pragma once
#include <Eigen/Dense>

class UpperBoundCalculatorInterface {
   public:
    virtual ~UpperBoundCalculatorInterface() = default;

// Запуск оптимизации. Возвращает максимальное значение функции
    virtual double calculate() = 0;

    void setNumThreads(int numThreads) {
        numThreads_ = numThreads;
    }
    void setLogInterval(int64_t interval) {
        logInterval_ = interval;
    }

// Возвращает точку куба, на которой достигнут максимум
    [[nodiscard]] virtual Eigen::RowVectorXd getMaxPoint() const = 0;

   protected:
    int     numThreads_  = 0;
    int64_t logInterval_ = 10'000'000;
};


=========================================================
// Файл: ./include/vertices_provider/HadamardFileProvider.h

#pragma once
#include "vertices_provider/ISimplexProvider.h"
#include "hadamard_utils/HadamardMatrixIterator.h"

#include <string>

class HadamardFileProvider : public ISimplexProvider {
   public:
    explicit HadamardFileProvider(const std::string& filename);
    std::vector<Eigen::RowVectorXd> getVertices() override;
    int                             getDimension() const override {
        return dim_;
    }
    HadamardMatrixIterator getIterator(const std::string& filename);

   private:
    static Eigen::MatrixXd readHadamardFile(const std::string& filename);
    Eigen::MatrixXd        H_;
    std::string            filename_;
    int                    dim_;
};

=========================================================
// Файл: ./include/vertices_provider/ISimplexProvider.h

#pragma once
#include <Eigen/Dense>
#include <vector>

class ISimplexProvider {
   public:
    virtual ~ISimplexProvider()                                  = default;
    virtual std::vector<Eigen::RowVectorXd> getVertices()        = 0;
    virtual int                             getDimension() const = 0;
};

=========================================================
// Файл: ./include/vertices_provider/SilvesterMethodProvider.h

#pragma once
#include "vertices_provider/ISimplexProvider.h"

class SilvesterMethodProvider : public ISimplexProvider {
   public:
    explicit SilvesterMethodProvider(int dim);
    std::vector<Eigen::RowVectorXd> getVertices() override;
    ;
    int getDimension() const override {
        return dim_;
    };
    static Eigen::MatrixXd hadamardMatrix(int order);

   private:
    int             dim_;
    Eigen::MatrixXd H_;
};

=========================================================
// Файл: ./include/vertices_provider/SimplexProviderFactory.h

#pragma once
#include "vertices_provider/ISimplexProvider.h"
#include <memory>
#include <string>

class SimplexProviderFactory {
   public:
    enum class Type { SILVESTER, FROM_HADAMARD_FILE };

    static std::unique_ptr<ISimplexProvider> create(Type type, int dim = 0,
                                                    const std::string& filename = "");
};

\end{lstlisting}
\end{document}